\section{Valid cuts from numerically chosen directions}\label{sec:certification}

A separation routine for~\eqref{eq:joint-hull} has two parts. One part
chooses a direction $(a,b)$; the other proves a lower bound on
$h_D(a,b)$. Only the second part affects validity. This section states
the resulting contract and the whole-domain bounds that we use for
non-quadratic expressions. Exact bounds for quadratic blocks follow in
Section~\ref{sec:quadratic}.

\subsection{The exported row}\label{sec:exported-row}

Floating-point LP solvers store binary64 coefficients. Every finite
binary64 number is a dyadic rational, so the row that the solver
actually uses is a rational inequality, and validity can be stated for
that row.

\begin{proposition}[Validity of the exported row]\label{prop:exported-row}
Let $\hat a\in\Q^n$ and $\hat b\in\Q^k$ be the exact values of the
coefficients passed to the solver. If
$L\le\min_{x\in D}\{\hat a\T x+\hat b\T F(x)\}$ and the exported
right-hand side satisfies $\hat r\le L$, then
$\hat a\T x+\hat b\T w\ge\hat r$ is valid for $H_D(F)$.
\end{proposition}

\begin{proof}
Apply Proposition~\ref{prop:support} to $(\hat a,\hat b)$ and weaken
the right-hand side.
\end{proof}

The statement is elementary, and the principle behind it is
established: a numerically chosen affine bound needs no verification of
how it was found, only a rigorous final shift
\citep[Section~5]{GarloffJanssonSmith2003JCAM}, \citep{GarloffSmith2008},
cut parameters may be chosen in floating point and the derivation then
repeated rigorously \citep[p.~294]{NeumaierShcherbina2004}, and safe
linear estimators are defined in terms of the machine coefficients
themselves \citep{BorradaileVanHentenryck2005}. What the proposition fixes
is what must be checked for joint support cuts. The direction may come
from an LP over samples, from a heuristic, or from a previous iteration;
none of these needs to be correct. A lower bound must be certified for the
\emph{final} coefficients, on the \emph{whole} domain $D$, and the
comparison $\hat r\le L$ must be exact. Rounding a proved rational row to
binary64 without accounting for the coefficient changes does not meet
this contract, and neither does a row that the solver changes after it
has been submitted: outside its exact mode, SCIP~10 rounds a row coefficient that is integral
within its tolerance to that integer, without adjusting the sides
(functions \code{rowAddCoef} and \code{rowChgCoefPos} in
\code{src/scip/lp.c} of \citealp{SCIPsource10}). For this reason our implementation reads back
the row that SCIP stores and compares it with the certified row
(Section~\ref{sec:replay}). When valid bounds $\ell\le(x,w)\le u$ are known
for all coordinates, there is a standard alternative: if $c\T z\ge L$ is
valid and $\hat c$ are the rounded coefficients, then
\begin{equation}\label{eq:bound-correction}
\hat c\T z\ \ge\ L+\sum_i\min\{(\hat c_i-c_i)\ell_i,\,(\hat c_i-c_i)u_i\}
\end{equation}
is valid, and its right-hand side can be rounded downward
\citep{NeumaierShcherbina2004,EiflerGleixner2024}. We use the direct
check of Proposition~\ref{prop:exported-row} for support inequalities
in lifted coordinates, because bounds on the lifted coordinates $w$
need not be available, and we use~\eqref{eq:bound-correction} after the
lifted coordinates have been eliminated (Section~\ref{sec:export}).

\subsection{Lower bounds on the whole domain}\label{sec:bounds}

\paragraph{Polynomials.}
Let $p$ be a polynomial with rational coefficients on a rational box
$B=\prod_{i=1}^n[\ell_i,u_i]$. After the affine change of variables
$x_i=\ell_i+(u_i-\ell_i)t_i$ and removal of fixed coordinates, write
\[
p(x(t))=\sum_{\beta\le d}a_\beta t^\beta=\sum_{\alpha\le d}c_\alpha B^d_\alpha(t),
\qquad
B^d_\alpha(t)=\prod_i\binom{d_i}{\alpha_i}t_i^{\alpha_i}(1-t_i)^{d_i-\alpha_i},
\]
where $d$ is a componentwise degree bound. The Bernstein coefficients
are
\begin{equation}\label{eq:bernstein}
c_\alpha=\sum_{\beta\le\alpha}a_\beta\prod_i
\binom{\alpha_i}{\beta_i}\Big/\binom{d_i}{\beta_i}.
\end{equation}

\begin{lemma}[Bernstein enclosure; see \citealp{GarloffSmith2008}]\label{lem:bernstein}
For every $x\in B$, $\min_\alpha c_\alpha\le p(x)\le\max_\alpha c_\alpha$.
\end{lemma}

\begin{proof}
On $[0,1]^n$ the basis polynomials are nonnegative and sum to one, so
$p(x(t))$ is a convex combination of the coefficients. Formula
\eqref{eq:bernstein} follows from the univariate identity
$t^j=\sum_{i\ge j}\binom ij\binom dj^{-1}B^d_i(t)$ applied in each
coordinate.
\end{proof}

All quantities in~\eqref{eq:bernstein} are rational, so the enclosure is
computed exactly. A weak enclosure is improved by subdivision, and the
enclosures shrink monotonically with the box
\citep{GarloffJanssonSmith2003}. Exact rational Bernstein subdivision has
even been formalized in a proof assistant \citep{MunozNarkawicz2013}; our
use of it is ordinary.

\begin{proposition}[Partition certificate]\label{prop:partition}
Let $B_1,\dots,B_N$ be closed boxes with $\bigcup_jB_j=B$, and let
$L_j\le\inf_{B_j}p$ for each $j$. Then $\min_jL_j\le\inf_Bp$. If a
rational affine row $\alpha\T x\le\gamma$ is part of the domain and
$\min_{x\in B_j}\alpha\T x>\gamma$, then $B_j$ may be omitted from the
minimum. If every cell is omitted, the domain is empty.
\end{proposition}

The proof is immediate. Its hypothesis is what a checker must verify:
the cells must cover $B$ exactly, starting from the trusted box and not
from a list supplied by the generator. A missing end interval or a
child box generated from a different parent would invalidate the bound.
For the same reason, a stored bound may be reused only for the same
expression, direction and domain; a bound proved on a smaller domain is
not valid on a larger one.

\paragraph{Univariate elementary expressions.}
For one variable we also admit sums, products and rational powers of
$\exp$, $\log$, $\sin$, $\cos$, $\sinh$, $\cosh$, $\tanh$, $\arctan$ and
$\abs{\cdot}$, with rational constants and $\pi$ and $e$. On every cell
of an interval subdivision, outward-rounded ball arithmetic
\citep{Johansson2017arb} gives an enclosure of the support objective
$t\mapsto at+b\T F(t)$. Domain requirements are part of the expression:
$\log$ needs a positive argument, a quotient needs a denominator interval
that excludes zero, and fractional powers use the nonnegative-base real
convention. If an enclosure is not finite or a domain requirement cannot
be proved on a cell, no bound is returned. Algebraic simplification
before this check is not allowed to remove a domain requirement; for
example, $x\log x/x$ still requires $x>0$.

A curvature bound often gives a much better lower bound on a cell than
a natural interval extension.
The following bound is the classical remainder of linear interpolation;
the same quantity is the maximal separation of an $\alpha$BB underestimator
with $\alpha=M/2$ \citep{AdjimanEtAl1998}.

\begin{lemma}[Chord correction]\label{lem:chord}
Let $p$ be twice continuously differentiable on $[\alpha,\beta]$ with
$p''\le M$. Then
\[
\min_{[\alpha,\beta]}p\ \ge\ \min\{p(\alpha),p(\beta)\}-\max\{0,M\}\,(\beta-\alpha)^2/8.
\]
\end{lemma}

\begin{proof}
The function $p(t)-Mt^2/2$ is concave, so it lies above its chord.
Rearranging gives
$p(t)\ge\ell_p(t)-\tfrac M2(t-\alpha)(\beta-t)$, where $\ell_p$ is the
chord of $p$. If $M\ge0$, use $(t-\alpha)(\beta-t)\le(\beta-\alpha)^2/4$
and $\ell_p(t)\ge\min\{p(\alpha),p(\beta)\}$. If $M<0$, the correction
term is nonnegative and $p\ge\ell_p$.
\end{proof}

Verified enclosures of $p(\alpha)$, $p(\beta)$ and of $p''$ on the cell
are enough. The implementation takes the larger of this bound and the
natural enclosure on each cell.

\paragraph{Outcomes.}
A bound computation returns one of three outcomes: a certified bound
with its partition, \emph{unsupported} (an operation outside the
grammar), or \emph{incomplete} (a work limit was reached). Only the
first produces a cut. An incomplete computation says nothing about the
point being separated; in particular it is not evidence that the point
lies in $H_D(F)$.

\subsection{A certificate for skipping separation}\label{sec:screening}

Separation is worth attempting only when a sufficiently violated cut
may exist. A convex combination of graph points certifies that no cut
with large normalized violation exists, and it is cheap to check.

Let $z=(x,w)\in\R^r$ collect the coordinates of the block, let $\bar z$
be the current LP point, and choose positive scales $s_1,\dots,s_r$.
Let $x^1,\dots,x^N\in D$ be samples whose graph values have verified
enclosures $z_j(x^i)\in[\ell_{ij},u_{ij}]$, and let
$\lambda\in\R^N_{\ge0}$ with $\sum_i\lambda_i=1$. Put
$L_j=\sum_i\lambda_i\ell_{ij}$, $U_j=\sum_i\lambda_iu_{ij}$, and
\[
\rho_j=\max\{\abs{\bar z_j-L_j},\abs{\bar z_j-U_j}\}/s_j,\qquad
R_1=\sum_j\rho_j,\qquad R_\infty=\max_j\rho_j.
\]

\begin{theorem}[Screening]\label{thm:screen}
Every valid inequality $c\T z\ge\beta$ for $H_D(F)$ satisfies
$\beta-c\T\bar z\le R_1$ if $\max_js_j\abs{c_j}\le1$, and
$\beta-c\T\bar z\le R_\infty$ if $\sum_js_j\abs{c_j}\le1$.
\end{theorem}

\begin{proof}
The point $z^\ast=\sum_i\lambda_iz(x^i)$ lies in $H_D(F)$ and satisfies
$L\le z^\ast\le U$. Hence
$\beta-c\T\bar z\le c\T(z^\ast-\bar z)\le\sum_j(s_j\abs{c_j})\,\rho_j$,
and H\"older's inequality gives both bounds.
\end{proof}

The theorem is the dual-norm bound on the violation of a hyperplane
\citep{Mangasarian1999} applied to a certified point of the hull; the
convex combination plays the role of an inclusion certificate in the
sense of \citet{Tawarmalani2010}.

The weights may come from any numerical LP. To use them, interpret the
returned values exactly, set negative entries to zero and divide by the
exact sum. Each sample must satisfy the exact domain constraints, and
the enclosures must come from the trusted evaluator rather than from the
LP. The direction LP of our implementation bounds each scaled
coefficient separately, so the matching certificate is $R_1$: if
$R_1\le\tau$, no cut with scaled violation above $\tau$ exists for that
block at this point, and the separation call can be skipped. A failed
screen has no consequence other than running the separator. In our first
campaign the screen was used by the automatic mode, and it allowed only
2 of 506 attempted skips; it is sound but rarely decisive, and the second
implementation does not use it.
